\chapter{Implementation}
\label{cha:implementation}

\epigraph{Ideas are easy, Implementation is hard.}{Guy Kawasaki}

Most of the implementation details and general decisions made for developing a
real-world functional project, based on the design architecture presented in chapter
\ref{cha:architecture}, are discussed in this chapter. It is organized in such a
way that it starts with the foundation and then progresses to the autoscaling of
worker nodes before moving on to example deployments of various applications in the
cluster. \\ %
The technology and programming languages utilized in the implementation are heterogeneous.
Nevertheless, according to a shared set of APIs and pre-defined structures, they
interact as a distinct and homogenous entity that maintains the entire cluster
in an active and healthy state. \\ %
Furthermore, several of the components and/or technologies discussed in this chapter
are interchangeable with other solutions. The latter is highly valuable for
organizations since it provides more configuration freedom and final control
over the cluster while continuously reflecting its ultimate goal.

\section{Distributions}
\label{sec:implementation_distributions}

% TODO Reference introduction dependencies

Almost all software relies on the necessity for a stable and consistent
operating system (OS). \\ %
The OS, as specified in section \ref{subsec:architecture_components_node}, is
installed on every node in the cluster and must be based on the \texttt{Linux
Kernel}. The latter is a critical requirement because most of the technologies
employed in the implementation rely on Linux primitives and core functionality
to operate properly. To enforce resource management for pods and containers, for
example, all Kubernetes distributions require the \texttt{cgroup v2}\footnote{\url{https://www.kernel.org/doc/html/latest/admin-guide/cgroup-v2.html}}
kernel module\footnote{\url{https://kubernetes.io/docs/concepts/architecture/cgroups}}.
Furthermore, the Linux Kernel is an Open Source project to which everyone can
contribute and get access, representing one fundamental objective of the entire project.
\\ %
To represent the potential of the implementation architecture being compatible with
several Linux distributions, the section is titled Distributions rather than
Operating System. A Linux distribution, often known as a Linux distro, is an operating
system that is built from components developed by various open-source projects and
programmers. The Linux kernel (the operating system's core), the GNU shell utilities
(terminal interface and commands), the X server (for a graphical desktop), the desktop
environment, a package management system, an installer, and other services are
all included in each distribution. Many components are created independently and
released in source code format. Thousands of software packages, utilities, and
applications can be found in a single Linux distribution. Linux distributions combine
code from open-source projects into a single operating system that can be installed
and booted. Linux distributions are available for a wide range of specialized
use cases, including desktop computers, servers without a graphical interface, supercomputers,
mobile devices, and more\cite{linux_distro}. \\ %
To be completely compatible with the implementation architecture, a Linux distribution
must meet three major requirements. They are described in the following list:
\begin{enumerate}
  \item The Linux Kernel version must be compatible with the Kubernetes version.
    \newline
    Newer versions of K8s may require newer versions of the Linux Kernel to
    provide newer functions and/or improve overall security.

  \item All necessary packages, as explained and documented in the section
    \ref{subsec:implementation_distributions_packages}, must be installed and globally
    accessible in the system.
    \newline
    A check on all essential packages is performed before the execution of the
    installation program, as detailed in section \ref{sec:implementation_installer}.
    If even one package is missing from the system, the installation fails with an
    error message identifying the missing packages.

  \item An Init System that is supported and compatible, as defined in the
    section \ref{subsec:implementation_distributions_init_system}.
    \newline
    OpenRC and Systemd are supported by the implementation. They are currently
    the two most used init systems, and the majority of available distributions are
    based on one of them.
\end{enumerate}
Finally, two custom Linux distributions are already configured, built, and
packaged as ISO files, as shown in section \ref{subsec:implementation_distributions_iso_image}.
They are based on a simple command line with no graphical interface and contain
only what is required to start a node. This is done to simplify the overall setup
by eliminating unnecessary software and technologies, addressing any
incompatibilities, and increasing the cluster's implementation and testing efficiency.
Both are derived from pre-existing Linux distributions that are known to be
stable, easy to configure, and lightweight, requiring just the bare minimum of resources
to operate. Furthermore, they are built on separate init systems, requiring the
overall implementation to be compatible with and tested with both.

\subsection{Packages}
\label{subsec:implementation_distributions_packages}

The packages used in the implementation that must be available in a Linux distribution
are specified in this section. A package is a collection of files that have been
bundled together and may be installed and uninstalled as a group\cite{package_manager}.
\\ %
All packages may be installed automatically using the package manager of the respective
Linux distribution. A package manager's job is to provide an interface that
assists the user in managing the collection of packages installed on the system.
A package manager maintains track of what software is installed on the system and
makes it simple to install new software, update software to newer versions, or delete
previously installed software. A package is usually a single piece of software. However,
it is typical for applications to be made up of numerous dependent packages. For
example, the \texttt{jq} package includes not just the \texttt{jq} command-line
JSON processor, but also the \texttt{oniguruma}\footnote{\url{https://github.com/kkos/oniguruma}}
package, a regular expressions library. Furthermore, various optional add-on packages
including extra data and documentation are available\cite{package_manager}. \\ %
All packages are already included by default in the two custom Linux
distributions. As a result, there is no need for any changes or the installation
of extra packages. \\ %
Each package in table \ref{tbl:packages} contains a brief description extracted
from its official documentation, as well as some notes on why it was utilized.
Furthermore, the required column indicates with symbols if the package is required
(\textcolor{bulmaGreen}{\faicon{check}}), optional (\textcolor{bulmaBlue}{\faicon{question}}),
or not required (\textcolor{bulmaRed}{\faicon{times}}). It should be noted that the
name of some packages may differ based on the package manager employed by the Linux
distribution.

\begin{xltabular}
  {\textwidth} { >{\ttfamily}l | c | X | X }

  \multicolumn{1}{ c |}{\large{\textbf{Name}}} &
  \multicolumn{1}{ c |}{\large{\textbf{Required}}} &
  \multicolumn{1}{ c |}{\large{\textbf{Description}}} &
  \multicolumn{1}{ c }{\large{\textbf{Notes}}} \\ \hline \hline

  bc & \textcolor{bulmaGreen}{\faicon{check}} & \texttt{bc}, also known as basic
  calculator, is a language that allows for arbitrary precision numbers with interactive
  statement execution\cite{bc}. & The Linux terminal does not support complicated
  math operations or number comparisons. As a result, the \texttt{bc} program is
  used throughout the installation procedure to guarantee that computed data, such
  as the node's mean power consumption and Simple Standard Deviation\footnote{\url{https://wikipedia.org/wiki/Standard_deviation}},
  is neither less than nor more than a predefined threshold. \\ \hline

  coreutils & \textcolor{bulmaGreen}{\faicon{check}} & The GNU Core Utilities are
  the GNU\footnote{\url{https://www.gnu.org}} operating system's essential file,
  shell, and text manipulation utilities. These are the core utilities that are
  required to be included in any operating system\cite{coreutils}. & By default,
  minimal Linux distributions only include a portion of the massive list of utilities
  provided in the standard \texttt{coreutils} package. As a result, several of
  the applications utilized in the implementation are missing. Therefore, the
  whole \texttt{coreutils} package is required. \\ \hline

  ethtool & \textcolor{bulmaGreen}{\faicon{check}} & \texttt{ethtool} is a standard
  Linux software for managing network drivers and hardware by modifying Network Interface
  Controller (NIC) settings\cite{ethtool}. & If the node's NIC interface supports
  Wake-on-Lan (WoL), \texttt{ethtool} is used to activate it.
  \newline
  It's worth noting that the NIC interface may reset to factory settings
  following a reboot. As a result, \texttt{ethtool} is always run before the NIC
  interface is activated. \\ \hline

  grep & \textcolor{bulmaGreen}{\faicon{check}} & \texttt{grep} prints lines from
  input files or standard input that match one or more patterns\cite{grep}. & \texttt{grep}
  is an essential utility in all Linux systems. It is widely utilized in all scripts
  for matching predetermined patterns and executing arbitrary actions based on the
  latter's results. \\ \hline

  inotify-tools & \textcolor{bulmaGreen}{\faicon{check}} & Monitor one or more
  files or directories for a predefined set of events, such as open, close, move/rename,
  delete or create. Under the hood, to be more efficient, it makes use of Linux's
  \texttt{inotify}\footnote{\url{https://docs.kernel.org/filesystems/inotify.html}}
  interface\cite{inotify_tools}. & Obtaining access information to the
  underlying Kubernetes cluster is required during the cluster initialization
  phase. This information is only available in an automatically generated configuration
  file saved in a specific and well-known directory. As a result, before proceeding
  to the next phase, the installation script uses the \texttt{inotifywait} program
  to wait for the file to be created. \\ \hline

  iproute2 & \textcolor{bulmaGreen}{\faicon{check}} & \texttt{iproute2} is a collection
  of utilities for managing TCP/IP\footnote{\url{https://wikipedia.org/wiki/Internet_protocol_suite}}
  networking and traffic control in Linux\cite{iproute2}.
  \newline
  The most essential tools are \texttt{ip} and \texttt{tc}. IPv4 and IPv6
  configuration is handled by \texttt{ip}, whereas \texttt{tc} stands for traffic
  control. & During the installation phase, the command \texttt{ip} is used to
  gather information about physical network interfaces that are not \texttt{loopback}\footnote{\url{https://wikipedia.org/wiki/Loopback}},
  such as name and MAC address. \\ \hline

  jq & \textcolor{bulmaGreen}{\faicon{check}} & \texttt{jq} is a lightweight and
  flexible command-line JSON\footnote{\url{https://www.json.org}} processor\cite{jq}.
  & To manipulate JSON data, all scripts make extensive use of \texttt{jq}. \\ \hline

  ncurses & \textcolor{bulmaRed}{\faicon{times}} & The \texttt{ncurses} (new curses)
  library is a free software emulation of curses in System V\footnote{\url{https://wikipedia.org/wiki/UNIX_System_V}}.
  It supports pads, color, multiple highlights, forms characters, and function-key
  mapping, as well as all of the other SVr4-curses improvements over BSD curses\cite{ncurses}.
  & The \texttt{ncurses} package includes the \texttt{tput} software, which is a
  utility for retrieving terminal capabilities in shell scripts.
  \newline
  When the spinner process is spawned in a script, \texttt{tput} is used to
  update its symbols at regular intervals, replacing the old one with the new one
  and updating the cursor position. Because a spinner is only a decorative
  element, the package is not required and can be omitted. \\ \hline

  nodejs & \textcolor{bulmaBlue}{\faicon{question}} & Node.js is an open-source,
  cross-platform JavaScript\footnote{\url{https://developer.mozilla.org/docs/Web/JavaScript}}
  runtime environment\cite{nodejs}. & This package is optional since the Server component,
  see section \ref{subsec:architecture_components_server}, may be bundled as a single
  binary or in a container, removing the need for the \texttt{nodejs} package to
  be installed. Because the present implementation does not provide any of the
  aforementioned alternatives, \texttt{nodejs} is included in both of the two
  custom Linux distributions. \\ \hline

  npm & \textcolor{bulmaBlue}{\faicon{question}} & The package manager for the Node.js
  JavaScript platform. It installs modules so that Node.js can locate them and
  intelligently resolves dependency conflicts. It is highly adaptable to serve a
  wide range of use cases. It is often used to publish, find, install, and build
  node applications\cite{npm}. & This package is optional since the Server component
  may be bundled as a single binary or in a container, removing the need for the
  \texttt{npm} package to be installed. Because the present implementation does
  not provide any of the aforementioned alternatives, \texttt{npm} is included
  in both of the two custom Linux distributions. \\ \hline

  openssh & \textcolor{bulmaGreen}{\faicon{check}} & The OpenSSH software suite includes
  tools for managing secure tunneling capabilities, various authentication
  methods, and advanced configuration options.
  \newline
  OpenSSH is the most popular remote login tool for the SSH protocol\footnote{\url{https://wikipedia.org/wiki/Secure_Shell}}.
  It encrypts all traffic to prevent eavesdropping, connection hijacking, and other
  types of attacks\cite{openssh}. & OpenSSH is critical in the cluster for protecting
  remote operations, key management, and other tasks. All operations are
  inherently unsafe without its software and can pose serious security issues, especially
  if the cluster is accessible from external networks. \\ \hline

  postgresql & \textcolor{bulmaBlue}{\faicon{question}} & \texttt{PostgreSQL} is
  a powerful, open-source object-relational database system with a high
  reputation for reliability, feature robustness, and performance\cite{postgresql}.
  & \texttt{PostgreSQL} is the database implementation used for the Database component
  (see section \ref{subsec:architecture_components_database}).
  \newline
  This package is optional since the database may be implemented outside of the
  cluster using other solutions. It is up to the organization in charge of the cluster
  to decide whether or not the database package should be included in the Linux distribution.
  \\ \hline

  procps-ng & \textcolor{bulmaRed}{\faicon{times}} & \texttt{procps} is a collection
  of command-line and full-screen programs that offer information from the \texttt{/proc}\footnote{\url{https://docs.kernel.org/filesystems/proc.html}}
  pseudo-filesystem. \texttt{procps} programs often concentrate on the structures
  that describe the processes operating on the system\cite{procps_ng}. & The \texttt{procps-ng}
  package includes the \texttt{ps} software, which displays information about
  the system's active processes.
  \newline
  When the spinner process is spawned in a script, \texttt{ps} is used to obtain
  the PID\footnote{\url{https://wikipedia.org/wiki/Process_identifier}} of the parent
  process. Because a spinner is only a decorative element, the package is not required
  and can be omitted. \\ \hline

  sed & \textcolor{bulmaGreen}{\faicon{check}} & \texttt{sed} is a stream editor.
  A stream editor is used on an input stream (file or standard input) to perform
  simple filtering and text transformations. While comparable to an editor that allows
  scripted edits, \texttt{sed} operates by making only one pass through the input(s)
  and is thus more efficient\cite{sed}. & \texttt{sed} is an essential utility in
  all Linux systems. It is widely utilized in all scripts for extracting well-defined
  portions and/or filtering data from its input(s). \\ \hline

  sudo & \textcolor{bulmaGreen}{\faicon{check}} & \texttt{sudo} (su\footnote{\url{https://wikipedia.org/wiki/Su_(Unix)}}
  "do") enables a system administrator to delegate authority to allow specific users
  (or groups of users) to run certain (or all) commands as root or another user while
  providing an audit record of the commands and their arguments. \texttt{sudo} works
  on a per-command basis; it is not a shell replacement\cite{sudo}. & The
  installation phase involves actions that need administrative authorization. Before
  performing any processing, the script checks to see whether the current user
  already has root\footnote{\url{https://wikipedia.org/wiki/Superuser}}
  privileges; if not, it requests them with \texttt{sudo}. \\ \hline

  sysbench & \textcolor{bulmaGreen}{\faicon{check}} & \texttt{sysbench} is a multithreaded
  scriptable benchmark tool for measuring CPU, memory, file I/O, mutex
  performance, and database performance\cite{sysbench}. & During the installation
  phase, \texttt{sysbench} is extensively used to analyze a cluster node's performance
  and, in conjunction with other tools, its power consumption. \\ \hline

  tar & \textcolor{bulmaGreen}{\faicon{check}} & \texttt{tar} is an archiving application
  that can store and manipulate numerous files in a single file (an archive).
  The archive can be either a standard file or a device, which can be placed on a
  local or remote system\cite{tar}. & \texttt{tar} is primarily used in
  production during the installation process to extract and manipulate various
  archive files.
  \newline
  \texttt{tar} is used in development to create the bundle archive for
  distribution, which contains all files and applications. \\ \hline

  tzdata & \textcolor{bulmaGreen}{\faicon{check}} & The Time Zone Database is a
  collection of code and data that represents the history of local time for
  several representative locations all over the world. It is updated regularly to
  reflect changes in time zone boundaries, UTC offsets, and daylight-saving laws
  enacted by legislative entities\cite{tzdata}. & To avoid date and time inconsistencies,
  all nodes in the cluster must be configured to the same timezone. By default,
  all nodes are configured to the Etc/UTC timezone. \\ \hline

  util-linux & \textcolor{bulmaGreen}{\faicon{check}} & \texttt{util-linux} is a
  random collection of Linux utilities\cite{util_linux}. & The installation
  script makes use of two \texttt{util-linux} utility programs: \texttt{lscpu} to
  display information about the CPU architecture and \texttt{lsblk} to list
  block devices. \\ \hline

  yq & \textcolor{bulmaGreen}{\faicon{check}} & \texttt{yq} is a portable and
  lightweight command-line YAML\footnote{\url{https://yaml.org}}, JSON, and XML\footnote{\url{https://www.w3.org/XML}}
  processor.
  \newline
  \texttt{yq} has a syntax similar to \texttt{jq}. It does not cover all that
  \texttt{jq} supports, but it does support the majority of the most common
  operations and functions\cite{yq}. & To manipulate YAML data, all scripts make
  extensive use of \texttt{yq}.
  \newline
  \texttt{yq} is primarily used to read and save YAML files. However, because of
  its more advanced functions and solutions, \texttt{jq} is used for the
  majority of complex data processing, at the expense of \texttt{yq}. As a result,
  YAML structures are first transformed to/from JSON, then processed with \texttt{jq},
  and then converted back to/from YAML. The double conversion should be
  eliminated and/or simplified in future implementations. \\ \hline

  \caption{Packages list}
  \label{tbl:packages}
\end{xltabular}

\subsection{Init System}
\label{subsec:implementation_distributions_init_system}

The Init System, short for Initialization System, is the first process that runs
when a Unix-based computer operating system is bootstrapped. Init is a daemon
process that is normally given Process IDentifier 1 (PID 1) and runs until the
system is shut down. It is the direct or indirect ancestor of all other processes
and adopts all orphaned processes automatically. The kernel starts Init during
the boot process; if the kernel is unable to start it, a kernel panic occurs\cite{init_system}.
\\ %
The Init System is logically positioned above the Kernel and is an essential component
of every modern system. It frequently integrates sophisticated features over
simple ones and is very customizable. Furthermore, it offers active control and
monitoring over the running processes, keeping the entire system healthy and active.
The latter is critical because it enables automated and continuous monitoring of
cluster implementation architecture-related processes operating in the node
without the need for a custom solution. For example, if a K8s instance executing
in a worker node crashes (i.e. exits with a code different than 0), the Init
System detects it and restarts the process automatically, making the node active
and able to receive workload again. \\ %
When the node is bootstrapped, the Init System is set to start all cluster-related
processes (i.e., K8s instance and Metric Server) and run a script that sends a
message to the server component notifying it that the node is active and ready. Before
the node is shut down, it automatically terminates all cluster-related processes
and sends a message to the server component informing that the node is shutting
down and will be inactive. As previously stated, the cluster performs these two
operations automatically: bootstrap for upscaling and shutdown for downscaling.
However, there are scenarios in which a node is manually bootstrapped and/or shut
down by human action without involving the server component directly. Despite
the possibility of the latter, the overall cluster state remains consistent, and
all node statuses are correctly updated, thanks to the Init System and its monitoring
and configuration capabilities. \\ %
The implementation currently supports two init systems: OpenRC (described in section
\ref{subsec:implementation_distributions_init_system_openrc}) and Systemd (described
in section \ref{subsec:implementation_distributions_init_system_systemd}). They are
the two most popular and widely used Init Systems, and the vast majority of
Linux distributions support one or both of them. Their capabilities are nearly identical.
However, configuration files and available application(s) that interact with the
bare-metal Init System API, are different, necessitating double effort to
achieve the same objective. \\ %
It should be noted that if the Init System daemon process crashes due to an
error, the entire system crashes and becomes inactive, necessitating a manual
restart. Even though the likelihood of the latter is extremely low, it is never
zero. As a result, as stated in section
\ref{subsec:architecture_components_server}, a heartbeat mechanism has been established
in the cluster to constantly monitor and identify the status and availability of
all active nodes. If there are no configuration or networking issues and a node's
heartbeat is not received within the predefined interval, it indicates that the heartbeat
daemon executing in the node has crashed and has not yet been restarted, or that
the entire node has crashed and is not automatically recoverable. \\ %
In each of the following sections, where an Init System distribution is described,
an example of the related service configuration file is presented to illustrate
its differences. A service configuration file contains information about a process
that the Init System controls and monitors. The application used as a reference
in this example is K3s, a lightweight Kubernetes distribution explained in further
detail in section \ref{subsec:implementation_dependencies_k3s}. Before starting,
it must clear any temporary files in the \texttt{tmp} directory and load
environment variables from two distinct directories. Furthermore, anytime the
application process crashes, the Init System must restart it within 5 seconds after
detection. Even though each Init System has a unique service configuration file,
the result is nearly identical. The list below highlights some of the most
significant configuration requirements. To aid comprehension, the identification
number for each requirement is also provided in the Init System's service
configuration file.
\begin{enumerate}[label=\protect\circled{\arabic{*}}]
  \item General information about the service.
    \newline
    Useful for users and/or administrators.

  \item Services on which the service is dependent.
    \newline
    The listed services must be activated before the current service may start.

  \item Remove any temporary files that begin with the name \texttt{k3s} and are
    placed in the temporary (\texttt{tmp}) directory before launching the application.

  \item Type of service. Different types correspond to various service behaviors
    as well as how the application is monitored and regulated.
    \newline
    In this case, monitor the application process and restart it if there is a
    crash or an anomaly.

  \item Absolute path to the application's binary file.

  \item Application arguments.

  \item Waiting time in seconds before restarting the crashed application.

  \item The maximum number of restarts permitted. When the threshold is met, the
    program is no longer restarted.
    \newline
    A value of 0 indicates that no threshold exists and that the application has
    no restart limit.

  \item Load \texttt{/etc/environment} and \texttt{/etc/rancher/k3s/k3s.env} environment
    variables files before starting the application.
\end{enumerate}

\subsubsection{OpenRC}
\label{subsec:implementation_distributions_init_system_openrc}

\begin{wrapfigure}
  {r}{.25\textwidth} %
  \centering
  \def\stackalignment{r}\stackunder{ \includegraphics[width=\linewidth]{images/logos/gentoo.png} } %
  {\scriptsize \parbox[t]{\linewidth}{ Source: \url{https://www.gentoo.org/inside-gentoo/artwork/gentoo-logo.html}} }
  \caption{Gentoo Linux logo}
\end{wrapfigure}

OpenRC is a dependency-based init system for Unix-like systems that maintain
compatibility with the system's init system, which is generally found in \texttt{/sbin/init}.
\\ %
At boot, OpenRC starts the appropriate system services in the right sequence, manages
them while the system is running, and stops them at shutdown. It can manage installed
daemons, optionally monitor the processes it launches, and start processes in parallel
when possible to reduce boot time. \\ %
OpenRC was developed for Gentoo\footnote{\url{https://www.gentoo.org}}, but it
may be used in other Linux distributions and BSD\footnote{\url{https://www.bsd.org}}
systems as well\cite{openrc}. \\ %
Listing \ref{lst:openrc} illustrates the example of the OpenRC service configuration
file for K3s saved in \texttt{/etc/init.d/k3s}. Manual control of the service is
possible using the \texttt{rc-service} command, which locates and runs the
OpenRC service with the specified arguments: \texttt{rc-service k3s <CMD> [...]}.
\texttt{<CMD>} can be substituted by \texttt{start} to start the service,
\texttt{stop} to stop the service, \texttt{restart} to restart the service, and
\texttt{status} to acquire the service status. Furthermore, with the command \texttt{rc-update
add k3s default}, the service may be set to start automatically when the system boots
up. The latter examples are merely a very small subset of all the available
commands/capabilities that OpenRC is capable of.

\begin{lstlisting}[language=shell, escapeinside={(*@}{@*)}, morekeywords={[2]{depend, start_pre, want, after, sourcex}}, morekeywords={[4]{name, description, supervisor, command, command_args, output_log, error_log, pidfile, respawn_delay, respawn_max}}, xleftmargin=\parindent, label={lst:openrc}, caption=\texttt{OpenRC} service configuration file for K3s]
  #!/sbin/openrc-run

  name="k3s"                                                                                      (*@ \circled{\texttt{\tiny{1}}} @*)
  description="Lightweight Kubernetes"                                                            (*@ \circled{\texttt{\tiny{1}}} @*)

  depend() {
    want cgroups                                                                                  (*@ \circled{\texttt{\tiny{2}}} @*)
    after network-online                                                                          (*@ \circled{\texttt{\tiny{2}}} @*)
  }

  start_pre() {
    rm -f "/tmp/k3s.*"                                                                            (*@ \circled{\texttt{\tiny{3}}} @*)
  }

  supervisor="supervise-daemon"                                                                   (*@ \circled{\texttt{\tiny{4}}} @*)
  command="/usr/local/bin/k3s"                                                                    (*@ \circled{\texttt{\tiny{5}}} @*)
  command_args="server"                                                                           (*@ \circled{\texttt{\tiny{6}}} @*)
  output_log="/var/log/k3s.log"
  error_log="/var/log/k3s.log"
  pidfile="/var/run/k3s.pid"
  respawn_delay=5                                                                                 (*@ \circled{\texttt{\tiny{7}}} @*)
  respawn_max=0                                                                                   (*@ \circled{\texttt{\tiny{8}}} @*)

  set -o allexport
  if [ -f "/etc/environment" ]; then sourcex "/etc/environment"; fi                               (*@ \circled{\texttt{\tiny{9}}} @*)
  if [ -f "/etc/rancher/k3s/k3s.env" ]; then sourcex "/etc/rancher/k3s/k3s.env"; fi               (*@ \circled{\texttt{\tiny{9}}} @*)
  set +o allexport
\end{lstlisting}

\subsubsection{systemd}
\label{subsec:implementation_distributions_init_system_systemd}

\begin{wrapfigure}
  {r}{.25\textwidth} %
  \centering
  \def\stackalignment{r}\stackunder{ \includegraphics[width=\linewidth]{images/logos/systemd.png} } %
  {\scriptsize \parbox[t]{\linewidth}{ Source: \url{https://brand.systemd.io}} }
  \caption{systemd logo}
\end{wrapfigure}

systemd is a collection of building blocks for a Linux system. It provides a
system and service manager that runs as PID 1 and starts the rest of the system.
\\ %
systemd supports aggressive parallelization, on-demand daemon startup, process
tracking through Linux control groups, mount and automount point management, and
a complex transactional dependency-based service control logic. \\ %
Other components include a logging daemon, tools to handle basic system configuration
such as the hostname, date, and locale, a list of logged-in users and running
containers and virtual machines, system accounts, runtime directories and
settings, and daemons to manage simple network configuration, network time
synchronization, log forwarding, and name resolution\cite{systemd}. \\ %
Listing \ref{lst:systemd} illustrates the example of the systemd service configuration
file for K3s saved in \texttt{/etc/systemd/system/k3s.service}. The \texttt{systemctl}
command, which manages the systemd system and the service manager, may be used to
manually control the service: \texttt{systemd <CMD> k3s [...]}. \texttt{<CMD>} can
be replaced with \texttt{start} to start the service, \texttt{stop} to stop it, \texttt{restart}
to restart it, and \texttt{status} to obtain the service status. Furthermore,
the service may be enabled to start automatically when the system boots up by
using the commands \texttt{systemctl enable k3s} followed by \texttt{systemctl daemon-reload}
to reload the systemd management configuration. The latter examples are only a small
portion of all the commands and capabilities provided by systemd.

\begin{lstlisting}[language=ini, escapeinside={(*@}{@*)}, morekeywords={Description, Documentation, Wants, After, StartLimitBurst, WantedBy, Type, ExecStartPre, ExecStart, EnvironmentFile, KillMode, Delegate, TimeoutStartSec, Restart, RestartSec}, xleftmargin=\parindent, label={lst:systemd}, caption=\texttt{systemd} service configuration file for K3s]
  [Unit]
  Description=Lightweight Kubernetes                                                              (*@ \circled{\texttt{\tiny{1}}} @*)
  Documentation=https://k3s.io                                                                    (*@ \circled{\texttt{\tiny{1}}} @*)
  Wants=network-online.target                                                                     (*@ \circled{\texttt{\tiny{2}}} @*)
  After=network-online.target                                                                     (*@ \circled{\texttt{\tiny{2}}} @*)
  StartLimitBurst=0                                                                               (*@ \circled{\texttt{\tiny{8}}} @*)

  [Install]
  WantedBy=multi-user.target

  [Service]
  Type=notify                                                                                     (*@ \circled{\texttt{\tiny{4}}} @*)
  ExecStartPre=/usr/bin/env sh -c 'rm -f "/tmp/k3s.*"'                                            (*@ \circled{\texttt{\tiny{3}}} @*)
  ExecStart=/usr/local/bin/k3s server                                                         (*@ \hspace*{1.375pt} \circled{\texttt{\tiny{5}}} \circled{\texttt{\tiny{6}}} @*)
  EnvironmentFile=-/etc/environment                                                               (*@ \circled{\texttt{\tiny{9}}} @*)
  EnvironmentFile=-/etc/rancher/k3s/k3s.env                                                       (*@ \circled{\texttt{\tiny{9}}} @*)
  KillMode=process
  Delegate=yes
  TimeoutStartSec=0
  Restart=always                                                                                  (*@ \circled{\texttt{\tiny{4}}} @*)
  RestartSec=5s                                                                                   (*@ \circled{\texttt{\tiny{7}}} @*)
\end{lstlisting}

\subsection{ISO image}
\label{subsec:implementation_distributions_iso_image}

A Linux distribution, including the two custom one, is packaged as a single ISO image
file. An optical disc image (or ISO image) is a disk image that contains
everything that would be written to an optical disc including the optical disc
file system. ISO images are expected to contain the binary image of an optical media
file system, including the data in its files in binary format, copied exactly as
they were stored on the disc. The data inside the ISO image will be structured
according to the file system that was used on the optical disc from which it was
created\cite{https://wikipedia.org/wiki/Optical_disc_image}.

The ISO image can be burned into physical media such as a CD, a USB flash drive
or can be mounted as an ISO file. The latter, shows the high level of
flexibility and overall compatibility that an ISO file has. Moreover, it is
extremely portable since it is intrinsecally only a single file making the
perfect format for a release bundle and very useful in an Air-Gap environment.

After a Linux Distribution's ISO image has been burned to the media it can be
installed on multiple systems with ease: attach the media to the system and select
it as default boot device. A live (volatile) system that is currently not
installed on the main disk starts. Then, the live Linux distribution can be
direclty installed into the main disk making it persistenlty available on the
system without the need of the removable media. The latter operation can be done
manually but it is more difficult than it sounds since requires knowledge on disk
formatting, partitioning and more. Fortunatly, most of the Linux distributions, including
the two custom ones, offers simple tools that automate the installation
procedure asking only a few simple questions.

In this section are briefly described the two original Linux distributions, respectively
Alpine Linux and Arch Linux, that has been used as the basis for the two custom
Linux distributions. Both of them have been used and tested toghether in the implementation
cluster. Moreover, it is shown how the two custom ISO image files are customized
and generated from a variaty of different configuration files merged together by
simple scripts.

\subsubsection{Generation}
\label{subsubsec:implementation_distributions_iso_generation}

\subsubsection{Alpine Linux}
\label{subsubsec:implementation_distributions_iso_alpine_linux}

\begin{wrapfigure}
  {r}{.25\textwidth} %
  \centering
  \def\stackalignment{r}\stackunder{ \includegraphics[width=\linewidth]{images/logos/alpine.png} } %
  {\scriptsize \parbox[t]{\linewidth}{ Source: \url{https://alpinelinux.org}} }
  \caption{Alpine Linux logo}
\end{wrapfigure}

Alpine Linux\footnote{\url{https://www.alpinelinux.org}} is a lightweight, security-focused,
and resource-efficient distribution. In recent years, has grown in popularity as
the foundation for the majority of Docker images. It is based on musl\footnote{\url{https://musl.libc.org}}
(an implementation of the C standard library, an alternative to glibc), BusyBox\footnote{\url{https://busybox.net}}
(a single, small executable that combines tiny versions of many common UNIX utilities),
a custom package manager called APK\footnote{\url{https://wiki.alpinelinux.org/wiki/Alpine_Package_Keeper}}
and the OpenRC init system ()\cite{alpine_linux}. \\ % TODO OpenRC reference
Alpine Linux is the preferred reCluster distribution. It has been utilized and
tested during development and is now present in every cluster node.

\subsubsection{Arch Linux}
\label{subsubsec:implementation_distributions_iso_arch_linux}

\begin{wrapfigure}
  {r}{.25\textwidth} %
  \centering
  \def\stackalignment{r}\stackunder{ \includegraphics[width=\linewidth]{images/logos/arch.png} } %
  {\scriptsize \parbox[t]{\linewidth}{ Source: \url{https://archlinux.org/art}} }
  \caption{Arch Linux logo}
\end{wrapfigure}

Arch Linux\footnote{\url{https://archlinux.org}} is a general-purpose
distribution that focuses on simplicity, minimalism, and code elegance. It is based
on Systemd init system () % TODO Systemd reference
and strives to stay bleeding edge, offering the latest stable versions of most
software through Pacman\footnote{\url{https://wiki.archlinux.org/title/Pacman}}
package manager. Uses a rolling release system that allows one-time installation
and perpetual software upgrades, allowing to not reinstall or upgrade the system
from one version to the next. Notably, Arch Linux is the foundation for several
many popular enterprise-grade distributions, such as Manjaro\footnote{\url{https://manjaro.org}}
and SteamOS\footnote{\url{https://store.steampowered.com/steamos}},
demonstrating its customizability, power, and stability\cite{arch_linux}.

\section{Dependencies}
\label{sec:implementation_dependencies}

% TODO Dependencies presentation

\subsection{Air-Gap Environment}
\label{subsec:implementation_dependencies_air_gap_environment}

\subsection{Management}
\label{subsec:implementation_dependencies_management}

\subsubsection{Configuration}
\label{subsec:implementation_dependencies_management_configuration}

\subsubsection{Script}
\label{subsec:implementation_dependencies_management_script}

\subsection{K3s}
\label{subsec:implementation_dependencies_k3s}

\begin{wrapfigure}
  {r}{.25\textwidth} %
  \centering
  \def\stackalignment{r}\stackunder{ \includegraphics[width=\linewidth]{images/logos/k3s.png} } %
  {\scriptsize \parbox[t]{\linewidth}{ Source: \url{https://k3s.io}} }
  \caption{K3s logo}
\end{wrapfigure}

\subsection{Node Exporter}
\label{subsec:implementation_dependencies_node_exporter}

\subsection{Prometheus}
\label{subsec:implementation_dependencies_prometheus}

\begin{wrapfigure}
  {r}{.25\textwidth} %
  \centering
  \def\stackalignment{r}\stackunder{ \includegraphics[width=\linewidth]{images/logos/prometheus.png} } %
  {\scriptsize \parbox[t]{\linewidth}{ Source: \url{https://prometheus.io}} }
  \caption{Prometheus logo}
\end{wrapfigure}

\subsection{Autoscaler}
\label{subsec:implementation_dependencies_autoscaler}

\subsubsection{Cluster Autoscaler}
\label{subsec:implementation_dependencies_autoscaler_cluster_autoscaler}

\section{Installer}
\label{sec:implementation_installer}

\subsection{POSIX Shell}
\label{subsec:implementation_installer_posix_shell}

\subsection{Configuration Arguments}
\label{subsec:implementation_installer_configuration_arguments}

\subsection{Configuration Files}
\label{subsec:implementation_installer_configuration_files}

\subsubsection{Certificates}
\label{subsubsec:implementation_installer_configuration_files_certificates}

\subsubsection{K3s}
\label{subsubsec:implementation_installer_configuration_files_k3s}

\subsubsection{K8s}
\label{subsubsec:implementation_installer_configuration_filesn_k8s}

\subsubsection{Node exporter}
\label{subsubsec:implementation_installer_configuration_files_node_exporter}

\subsubsection{reCluster}
\label{subsubsec:implementation_installer_configuration_files_recluster}

\subsubsection{SSH}
\label{subsubsec:implementation_installer_configuration_files_ssh}

\subsection{Node Facts}
\label{subsec:implementation_installer_node_facts}

\subsection{Node Information}
\label{subsec:implementation_installer_node_information}

\subsection{Node Benchmarks}
\label{subsec:implementation_installer_node_benchmarks}

\subsubsection{SysBench}
\label{subsubsec:implementation_installer_node_benchmarks_sysbench}

\subsection{Node Power Consumption}
\label{subsec:implementation_installer_node_power_consumption}

\subsubsection{CloudFree Smart Plug}
\label{subsubsec:implementation_installer_node_power_consumption_cloudfree_smart_plug}

\subsection{Node Registration}
\label{subsec:implementation_installer_node_registration}

\subsubsection{K8s Node Label And Name}
\label{subsubsec:implementation_installer_node_registration_k8s_node_label_and_name}

\subsection{Cluster Initialization}
\label{subsec:implementation_installer_cluster_initialization}

\subsubsection{Kubeconfig}
\label{subsubsec:implementation_installer_cluster_initialization_kubeconfig}

\subsubsection{Database}
\label{subsubsec:implementation_installer_cluster_initialization_database}

\subsubsection{Server}
\label{subsubsec:implementation_installer_cluster_initialization_server}

\subsubsection{Admin And Autoscaler Users}
\label{subsubsec:implementation_installer_cluster_initialization_admin_and_autoscaler_users}

\subsubsection{K8s}
\label{subsubsec:implementation_installer_cluster_initialization_k8s}

\subsection{Services}
\label{subsec:implementation_installer_services}

\subsubsection{Start}
\label{subsubsec:implementation_installer_services_start}

\subsubsection{Stop}
\label{subsubsec:implementation_installer_services_stop}

\subsubsection{Boot}
\label{subsubsec:implementation_installer_services_boot}

\subsubsection{Shutdown}
\label{subsubsec:implementation_installer_services_shutdown}

\section{Server}
\label{sec:implementation_server}

\subsection{Database}
\label{subsec:implementation_server_database}

\subsubsection{PostgreSQL}
\label{subsubsec:implementation_server_database_postgresql}

\begin{wrapfigure}
  {r}{.25\textwidth} %
  \centering
  \def\stackalignment{r}\stackunder{ \includegraphics[width=\linewidth]{images/logos/postgresql.png} } %
  {\scriptsize \parbox[t]{\linewidth}{ Source: \url{https://wiki.postgresql.org/wiki/Logo}} }
  \caption{PostgreSQL logo}
\end{wrapfigure}

\subsubsection{Prisma}
\label{subsubsec:implementation_server_database_prisma}

\begin{wrapfigure}
  {r}{.25\textwidth} %
  \centering
  \def\stackalignment{r}\stackunder{ \includegraphics[width=\linewidth]{images/logos/prisma.png} } %
  {\scriptsize \parbox[t]{\linewidth}{ Source: \url{https://www.prisma.io}} }
  \caption{Prisma logo}
\end{wrapfigure}

\subsubsection{Schema}
\label{subsubsec:implementation_server_database_schema}

\subsection{GraphQL}
\label{subsec:implementation_server_graphql}

\subsubsection{GraphQL vs REST}
\label{subsubsec:implementation_server_graphql_graphql_vs_rest}

\subsubsection{Schema}
\label{subsubsec:implementation_server_graphql_schema}

\subsection{Node Pool}
\label{subsec:implementation_server_node_pool}

\subsection{Node Registration}
\label{subsec:implementation_server_node_registration}

\subsubsection{JSON Web Token}
\label{subsubsec:implementation_server_node_registration_json_web_token}

\subsection{Upscaling}
\label{subsec:implementation_server_upscaling}

\subsubsection{Wake-on-LAN}
\label{subsubsec:implementation_server_scale_up_wake_on_lan}

\subsection{Downscaling}
\label{subsec:implementation_server_downscaling}

\subsubsection{SSH}
\label{subsubsec:implementation_server_scale_up_ssh}

\section{Autoscaling}
\label{sec:implementation_autoscaling}

\subsection{Vertical Pod Autoscaler}
\label{subsec:implementation_autoscaling_vertical_pod_autoscaler}

\subsection{Horizontal Pod Autoscaler}
\label{subsec:implementation_autoscaling_horizontal_pod_autoscaler}

\subsection{Cluster Autoscaler}
\label{subsec:implementation_autoscaling_cluster_autoscaler}

\subsubsection{Cloud Provider}
\label{subsubsec:implementation_autoscaling_cluster_autoscaler_cloud_provider}

\subsubsection{Configuration}
\label{subsubsec:implementation_autoscaling_cluster_autoscaler_configuration}

% TODO Cloud Configuration

\subsubsection{Upscaling}
\label{subsubsec:implementation_autoscaling_cluster_autoscaler_upscaling}

\subsubsection{Downscaling}
\label{subsubsec:implementation_autoscaling_cluster_autoscaler_downscaling}